\chapter*{Как пользоваться сборником}
\addcontentsline{toc}{chapter}{Как пользоваться сборником}
Основа --- программа-минимум физфака МГУ 2024 года~\cite{program2024}: шесть разделов, 49 рабочих вопросов и 100 подвопросов. Это организация подготовки, а не официальный комплект билетов. Заголовки соответствуют программе; полная формулировка приведена у группы, модели и примеры --- внутри ответа.

\paragraph{Номера и переходы.}
Первая цифра --- раздел, вторая --- вопрос, третья --- самостоятельный подвопрос. В оглавлении только шесть тематических разделов и служебные страницы. Полное деление находится в \textbf{списке вопросов}: название ведёт к тексту, стрелка $\uparrow$ у заголовка --- обратно к конкретной строке. Подвопрос можно готовить отдельно, переходя по ссылкам к используемым определениям и выводам.

\paragraph{Порядок ответа.}
Внутри ответа: определения и постановка, физический механизм, основные уравнения и переходы, смысл, пример, границы применимости. Рисунки помогают объяснению на доске. Заключительные подсказки обычно общие для вопроса:
\begin{itemize}
\item \textbf{Короткий ответ / для повторения} --- связный маршрут рассказа.
\item \textbf{На доске} --- основные уравнения, графики или схема.
\item \textbf{Вопрос для себя} --- проверка понимания.
\item \textbf{Удобный вопрос} --- уже разобранная деталь для возможного продолжения.
\item \textbf{Источники} --- использованные места; номер ссылки ведёт к библиографии с URL.
\end{itemize}
\textbf{Зачем «удобный вопрос».} Основную цепочку устного ответа рассказывают целиком, а второстепенную деталь можно намеренно не разворачивать, оставив повод для уточнения. Ответ уже есть выше: например, смысл предельного случая. Это возможность направить разговор к знакомому материалу; комиссия может спросить другое. Следовать подсказке или сразу рассказать всё --- личный выбор.

\paragraph{Условные обозначения.}
\begin{center}\small
\begin{tabular}{cl}
\toprule Знак & Значение\\\midrule
\flagicon & Только заготовка и разметка\\
\clockicon & Частичный текст, находится в работе\\
\bookicon & Полный текст, пригоден для чтения\\
В.\reviewunread\quad Д.\reviewunread & Соответствующий человек не проверял эту версию\\
В.\reviewread\quad Д.\reviewread & Прочитано, окончательно не принято\\
В.\reviewaccepted\quad Д.\reviewaccepted & Принято соответствующим рецензентом\\\bottomrule
\end{tabular}\end{center}
В. и Д. --- два независимых рецензента. Цвет продублирован формой. Книга и успешная компиляция не означают человеческого принятия. Существенно изменённый текст требует повторного просмотра; все ответы этой редакции ещё ожидают принятия. Состояние группы учитывает её общий материал и подвопросы.

